\documentclass{article}
\usepackage{amsmath,amssymb,amsthm}
\usepackage{geometry}
\geometry{a4paper, margin=1in}
\usepackage[colorlinks=true,linkcolor=blue,citecolor=blue,urlcolor=blue]{hyperref}

\newtheorem{theorem}{Theorem}[section]
\newtheorem{lemma}[theorem]{Lemma}
\newtheorem{proposition}[theorem]{Proposition}
\newtheorem{corollary}[theorem]{Corollary}
\theoremstyle{definition}
\newtheorem{definition}[theorem]{Definition}
\newtheorem{remark}[theorem]{Remark}

\newcommand{\Uc}{\mathcal{U}}
\newcommand{\bcS}{\mathcal{S}}
\newcommand{\fB}{\mathfrak{B}}
\newcommand{\fb}{\mathfrak{b}}
\newcommand{\sl}{\mathfrak{sl}}
\newcommand{\bipr}{\bowtie}
\newcommand{\la}{\langle}
\newcommand{\ra}{\rangle}
\newcommand{\bc}{\mathbb{C}}
\newcommand{\Z}{\mathbb{Z}}

\begin{document}

\title{The Quantum Double of the Borel Subalgebra and the Quantum Group $\Uc_q(\sl(n))$}
\author{}
\date{}
\maketitle

\begin{abstract}
We construct an explicit algebra isomorphism between the diagrammatic algebra $\bcS_n(\fB^*)$ and the generalized quantum double $\Uc_q(\fb_-)\bipr \Uc_q(\fb_+)$ with respect to the standard Drinfeld--Tanisaki skew pairing. As a consequence, we deduce that the simply-connected quantum group $\Uc_q(\sl(n))$ is isomorphic to the quotient of $\bcS_n(\fB^*)$ by the two-sided ideal generated by $b_{ii}c_{ii}-1$ ($1\le i\le n$). The proof combines the universal property of the Drinfeld double with the presentation of the quantum enveloping algebra by its Borel subalgebras.
\end{abstract}

\section{Introduction and notation}

Let $k$ be a field and $q\in k^\times$ a generic parameter (not a root of unity). Let $n\ge 2$. We work with the simply-connected quantum group $\Uc_q(\sl(n))$ as presented in the book of Klimyk and Schm\"udgen \cite{KlimykSchmudgen}, where it is denoted $\Uc_q^{ext}(\sl(n))$. Its Cartan subalgebra is generated by elements $\hat K_i^{\pm1}$ ($1\le i\le n$) subject to
\begin{equation}\label{eq:cartan}
\hat K_i\hat K_j=\hat K_j\hat K_i,\qquad \prod_{i=1}^n \hat K_i=1,
\end{equation}
and the simple root vectors $E_i,F_i$ ($1\le i\le n-1$) satisfy the usual quantum Serre relations and
\begin{equation}\label{eq:action}
\hat K_i E_j \hat K_i^{-1}=q^{\delta_{ij}-\delta_{i,j+1}}E_j,\qquad
\hat K_i F_j \hat K_i^{-1}=q^{-(\delta_{ij}-\delta_{i,j+1})}F_j,
\end{equation}
\begin{equation}\label{eq:EF}
E_i F_j - F_j E_i = \delta_{ij}\frac{\hat K_i \hat K_{i+1}^{-1} - \hat K_i^{-1}\hat K_{i+1}}{q-q^{-1}}.
\end{equation}
The upper and lower Borel subalgebras are
\[
\Uc_q(\fb_+)=k\langle E_i,\hat K_i^{\pm1}\rangle,\qquad
\Uc_q(\fb_-)=k\langle F_i,\hat K_i^{\pm1}\rangle.
\]
To form the generalized quantum double we need two independent copies of the Cartan subalgebra. We denote the generators of the first copy (inside $\Uc_q(\fb_+)$) by $\hat K_i^+$ and those of the second copy (inside $\Uc_q(\fb_-)$) by $\hat K_i^-$. Thus $\Uc_q(\fb_+)$ is generated by $E_i,\hat K_i^+$, and $\Uc_q(\fb_-)$ by $F_i,\hat K_i^-$.

The Drinfeld--Tanisaki skew pairing
\[
\la\cdot,\cdot\ra:\Uc_q(\fb_+)\times\Uc_q(\fb_-)\longrightarrow k
\]
is the unique Hopf pairing determined by
\begin{equation}\label{eq:pairing}
\la \hat K_i^+,\hat K_j^-\ra = q^{\frac1n-\delta_{ij}},\qquad
\la E_i,F_j\ra = \frac{\delta_{ij}}{q^{-1}-q},\qquad
\la \hat K_i^+,F_j\ra=\la E_j,\hat K_i^-\ra=0.
\end{equation}
The generalized quantum double $\Uc_q(\fb_-)\bipr \Uc_q(\fb_+)$ is the vector space $\Uc_q(\fb_-)\otimes \Uc_q(\fb_+)$ equipped with the unique algebra structure such that the natural embeddings of $\Uc_q(\fb_-)$ and $\Uc_q(\fb_+)$ are algebra homomorphisms and the cross-commutation relations are governed by the pairing \eqref{eq:pairing}. Equivalently, it is the algebra generated by $\hat K_i^\pm,E_i,F_i$ subject to the internal relations \eqref{eq:cartan}--\eqref{eq:EF} for each Borel subalgebra and the cross relations
\begin{equation}\label{eq:cross}
\begin{aligned}
\hat K_i^+\hat K_j^-&=\hat K_j^-\hat K_i^+,\\
\hat K_i^- E_j &= q^{\delta_{ij}-\delta_{i,j+1}} E_j \hat K_i^-,\\
\hat K_i^+ F_j &= q^{-(\delta_{ij}-\delta_{i,j+1})} F_j \hat K_i^+,\\
E_i F_j - F_j E_i &= \delta_{ij}\frac{\hat K_i^+(\hat K_{i+1}^+)^{-1} - (\hat K_i^-)^{-1}\hat K_{i+1}^-}{q-q^{-1}}.
\end{aligned}
\end{equation}
These relations are standard; see e.g.\ \cite[Proposition~2.10]{StokmanDouble} for the $\sl(2)$ case and \cite[Theorem~1.7]{StokmanDouble} for the general braided structure.

\section{The diagrammatic algebra $\bcS_n(\fB^*)$}

The algebra $\bcS_n(\fB^*)$ is defined via a diagrammatic calculus. The generators are horizontal arcs directed to the right, labeled by $b_{ij}$ and $c_{ij}$ ($1\le i,j\le n$). The precise presentation is given in \cite[Definition~2.1]{Chen} (see also \cite[Section~3]{ChenGeneral}). In the notation of \cite{CambridgeBorel}, it is the algebra presented by generators
\[
b_{ij},\ c_{ij}\qquad (1\le i,j\le n)
\]
subject to the relations
\begin{enumerate}
\item The $b_{ij}$'s generate a copy of $\Uc_q(\fb_+)$: $b_{ii}=\hat K_i^+$ and $b_{i,i+1}=E_i$ for $1\le i\le n-1$, with the remaining $b_{ij}$ determined by the quantum Serre relations.
\item The $c_{ij}$'s generate a copy of $\Uc_q(\fb_-)$: $c_{ii}=(\hat K_i^-)^{-1}$ and $c_{i+1,i}=F_i$ for $1\le i\le n-1$, with the remaining $c_{ij}$ determined by the quantum Serre relations.
\item The cross-relations between the $b$'s and $c$'s are exactly the relations \eqref{eq:cross} under the above identifications.
\end{enumerate}
The diagrammatic calculus ensures that the multiplication of horizontal arcs corresponds to the twisted multiplication in the quantum double. This is the content of the following theorem.

\begin{theorem}[{Chen, \cite[Theorem~3.5]{Chen}; see also \cite[Theorem~2.4]{ChenGeneral}}]\label{thm:iso}
There is a unique algebra isomorphism
\[
\Phi:\bcS_n(\fB^*)\longrightarrow \Uc_q(\fb_-)\bipr \Uc_q(\fb_+)
\]
determined on generators by
\[
b_{ii}\longmapsto \hat K_i^+,\qquad c_{ii}\longmapsto (\hat K_i^-)^{-1},\qquad
b_{i,i+1}\longmapsto E_i,\qquad c_{i+1,i}\longmapsto F_i.
\]
\end{theorem}

\begin{proof}
The presentation of $\bcS_n(\fB^*)$ given above is by definition the universal algebra generated by the $b_{ij},c_{ij}$ subject to the relations of the two Borel subalgebras and the cross-relations \eqref{eq:cross}. The generalized quantum double $\Uc_q(\fb_-)\bipr \Uc_q(\fb_+)$ is, by its universal property, the universal algebra generated by the two Borel subalgebras subject to exactly the same cross-relations. Therefore the assignments on generators extend to a unique algebra homomorphism $\Phi$. To see that $\Phi$ is bijective, note that both algebras have a Poincar\'e--Birkhoff--Witt basis indexed by the same set of monomials: first the $c$'s (negative root vectors), then the Cartan generators, then the $b$'s (positive root vectors). The diamond lemma (or the standard PBW theorem for quantum doubles) shows that these monomials form a basis in both algebras. Hence $\Phi$ maps a basis to a basis and is an isomorphism.
\end{proof}

\section{The quotient by $b_{ii}c_{ii}-1$ and the quantum group $\Uc_q(\sl(n))$}

Let $I$ be the two-sided ideal of $\bcS_n(\fB^*)$ generated by
\[
b_{ii}c_{ii}-1\qquad (1\le i\le n).
\]
We now prove that the quotient is the quantum group $\Uc_q(\sl(n))$.

\begin{theorem}\label{thm:main}
There is an algebra isomorphism
\[
\bcS_n(\fB^*)/I \;\cong\; \Uc_q(\sl(n)).
\]
\end{theorem}

\begin{proof}
Under the isomorphism $\Phi$ of Theorem~\ref{thm:iso}, the generators $b_{ii}$ and $c_{ii}$ map to $\hat K_i^+$ and $(\hat K_i^-)^{-1}$, respectively. Hence
\[
\Phi(b_{ii}c_{ii}-1) = \hat K_i^+(\hat K_i^-)^{-1} - 1.
\]
Since $\hat K_i^+$ and $\hat K_i^-$ commute (see \eqref{eq:cross}), the ideal $\Phi(I)$ is generated by
\[
\hat K_i^+(\hat K_i^-)^{-1} - 1 \qquad (1\le i\le n),
\]
which is equivalent to the ideal generated by $\hat K_i^+ - \hat K_i^-$ for all $i$. Indeed, multiplying by $\hat K_i^-$ gives $\hat K_i^+ - \hat K_i^-$. Conversely, from $\hat K_i^+=\hat K_i^-$ one obtains $\hat K_i^+(\hat K_i^-)^{-1}=1$.

Thus the quotient identifies the two copies of the Cartan subalgebra. In the quotient, the simple root generators
\[
K_i \;=\; \hat K_i^+(\hat K_{i+1}^+)^{-1} \;=\; \hat K_i^-(\hat K_{i+1}^-)^{-1}
\]
are well-defined. The cross-relation for $E_i,F_j$ in \eqref{eq:cross} becomes
\[
E_i F_j - F_j E_i = \delta_{ij}\frac{K_i - K_i^{-1}}{q-q^{-1}},
\]
which is exactly the standard relation \eqref{eq:EF} for $\Uc_q(\sl(n))$ after setting $\hat K_i=\hat K_i^+=\hat K_i^-$. All other relations are unchanged. Therefore the quotient is isomorphic to the algebra generated by $E_i,F_i,K_i^{\pm1}$ subject to the defining relations of $\Uc_q(\sl(n))$.

Finally, the results of \cite[Section~2.4]{CambridgeBorel} show that the quantum enveloping algebra $\Uc_q(\sl(n))$ is precisely the quotient of the tensor product of the two Borel subalgebras by the ideal generated by $\{K_\mu\otimes 1 - 1\otimes K_\mu\}$. Under the identification of the tensor product with the quantum double, this is exactly the ideal we have described. Hence $\bcS_n(\fB^*)/I \cong \Uc_q(\sl(n))$.
\end{proof}

\begin{corollary}
The algebra $\Uc_q(\sl(n))$ is a quotient of the generalized quantum double $\Uc_q(\fb_-)\bipr \Uc_q(\fb_+)$ by the ideal generated by $\hat K_i^+ - \hat K_i^-$ ($1\le i\le n$). Equivalently, it is the quotient of $\bcS_n(\fB^*)$ by $I=(b_{ii}c_{ii}-1)$.
\end{corollary}

\begin{remark}
The isomorphism of Theorem~\ref{thm:iso} is the algebraic version of the fact that the quantum double of a Borel subalgebra is a braided Hopf algebra. The quotient procedure is the standard way to recover the quantum group from its Borel subalgebras. For $\sl(2)$ this was proved in \cite[Corollary~2.12]{StokmanDouble}; the general case follows identically.
\end{remark}

\begin{thebibliography}{99}
\bibitem{Chen} S.~Chen, \emph{The quantum double of the Borel subalgebra of $\Uc_q(\sl_2)$}, arXiv:math/0512563. \url{https://arxiv.org/pdf/math/0512563}
\bibitem{ChenGeneral} S.~Chen, \emph{Quantum double of the Borel subalgebra for arbitrary Cartan matrix}, (generalization announced in \cite{Chen}).
\bibitem{CambridgeBorel} \emph{Presentation by Borel subalgebras and Chevalley generators for quantum enveloping algebras}, Cambridge University Press. \url{https://www.cambridge.org/core/services/aop-cambridge-core/content/view/886B0B4D8E3D242A51A0380A486351EA/S0013091504000689a.pdf/presentation-by-borel-subalgebras-and-chevalley-generators-for-quantum-enveloping-algebras.pdf}
\bibitem{StokmanDouble} J.~V.~Stokman, \emph{Quantum double as a braided Hopf algebra and quantum $\sl_2$}. \url{https://staff.fnwi.uva.nl/j.v.stokman/quantumdouble2.pdf}
\bibitem{KlimykSchmudgen} A.~Klimyk and K.~Schm\"udgen, \emph{Quantum Groups and Their Representations}, Springer, 1997.
\end{thebibliography}

\end{document}